\section{Spectral barriers}\label{sec:barriers}

\subsection{An abstract barrier lemma}

In this subsection $Q=V+K$ is a bounded self-adjoint operator on $\ell^2(\Z)$, $V$ is a real diagonal operator, $K$ has range at most $q$ (that is, $\langle\delta_i,K\delta_j\rangle=0$ for $|i-j|>q$), and $\|K\|\le J$. Fix $d_0\in\R$ and $W>0$, and let $\Pi=\un_{[d_0-W,d_0+W]}(Q)$. For an interval $I\subset\Z$ write $\partial_qI=\{i\in I:\dist(i,\Z\setminus I)\le q\}$.

\begin{lemma}[Outer barrier]\label{lem:outer}
Let $I\subset\Z$ be an interval on which either $V-d_0\ge V_0$ or $V-d_0\le-V_0$, where $V_0\ge W+2J$ and $V_0>W$. If $B\subset I$ and $\dist(B,\Z\setminus I)>(\ell+1)q$ for an integer $\ell\ge0$, then
\[
\|\un_B\Pi\|\le2\Bigl(\frac{J}{V_0-W}\Bigr)^{\ell+2}.
\]
\end{lemma}

\begin{proof}
Assume $V-d_0\ge V_0$ on $I$; the other case follows by replacing $Q$ and $d_0$ with $-Q$ and $-d_0$. Let $Q_I=\un_IQ\un_I$ acting on $\ell^2(I)$, let $A$ be the restriction of $Q$ to $\operatorname{Ran}\Pi$, and let $Y=-\un_IQ\un_{\Z\setminus I}\Pi$. Since $\Pi$ commutes with $Q$, the operator $X=\un_I\Pi$ solves the Sylvester equation
\[
(Q_I-d_0)X-X(A-d_0)=Y.
\]
Here $Q_I-d_0\ge V_0-J$ and $\|A-d_0\|\le W<V_0-J$; the strict inequality follows from $V_0\ge W+2J$ if $J>0$, and from $V_0>W$ if $J=0$. So
\[
X=\int_0^\infty e^{-s(Q_I-d_0)}\,Y\,e^{s(A-d_0)}\,ds,
\]
as one checks by differentiating the integrand. Since $V$ is diagonal, $Y=-\un_IK\un_{\Z\setminus I}\Pi$, so $\|Y\|\le J$ and the range of $Y$ lies in $\ell^2(\partial_qI)$.

Expand $e^{-s(Q_I-d_0)}$ in a Dyson series around the diagonal part $V_I-d_0\ge V_0$. The term with $k$ factors of $K_I=\un_IK\un_I$ has norm at most $e^{-sV_0}(sJ)^k/k!$ and moves supports by at most $kq$. Since $\dist(B,\partial_qI)>\ell q$, only terms with $k\ge\ell+1$ connect $\ell^2(\partial_qI)$ to $\ell^2(B)$. Therefore
\[
\|\un_BX\|\le J\int_0^\infty e^{-s(V_0-W)}\sum_{k\ge\ell+1}\frac{(sJ)^k}{k!}\,ds=\sum_{k\ge\ell+1}\Bigl(\frac{J}{V_0-W}\Bigr)^{k+1}\le2\Bigl(\frac{J}{V_0-W}\Bigr)^{\ell+2},
\]
using $J/(V_0-W)\le\frac12$.
\end{proof}

\begin{lemma}[Inner core]\label{lem:inner}
Let $I_0\subset\Z$ be an interval with $|V-d_0|\le W/4$ on $I_0$, and assume $W\ge4J$. If $0\in I_0$ and $\dist(0,\Z\setminus I_0)>(\ell+1)q$, then
\[
\|(1-\Pi)\delta_0\|\le3\Bigl(\frac{4J}{3W}\Bigr)^{\ell+2}.
\]
\end{lemma}

\begin{proof}
Let $\Pi_+=\un_{(d_0+W,\infty)}(Q)$ and let $A_+$ be the restriction of $Q$ to its range, so $A_+-d_0\ge W$. With $Q_{I_0}=\un_{I_0}Q\un_{I_0}$ we have $Q_{I_0}-d_0\le W/4+J$, and $X_+=\un_{I_0}\Pi_+$ solves $(Q_{I_0}-d_0)X_+-X_+(A_+-d_0)=Y_+$ with $Y_+=-\un_{I_0}K\un_{\Z\setminus I_0}\Pi_+$. Hence
\[
X_+=-\int_0^\infty e^{s(Q_{I_0}-d_0)}\,Y_+\,e^{-s(A_+-d_0)}\,ds.
\]
Expanding $e^{s(Q_{I_0}-d_0)}$ around its diagonal part, which has norm at most $W/4$, only terms with at least $\ell+1$ factors of $K$ connect $\partial_qI_0$ to the site $0$. As in Lemma~\ref{lem:outer},
\[
\|\Pi_+\delta_0\|=\|\delta_0^*X_+\|\le J\int_0^\infty e^{-3sW/4}\sum_{k\ge\ell+1}\frac{(sJ)^k}{k!}\,ds\le2\Bigl(\frac{4J}{3W}\Bigr)^{\ell+2}.
\]
The same bound holds for $\Pi_-=\un_{(-\infty,d_0-W)}(Q)$, applying the argument to $-Q$. Since $1-\Pi=\Pi_++\Pi_-$ with orthogonal ranges, the claim follows.
\end{proof}

\begin{lemma}[Cutoff]\label{lem:cutoff}
Let $\Pi$ be a spectral projection of an operator that commutes with $H$, let $N\ge1$ be an integer, and put $B_N=\{-N-1,-N,N,N+1\}$. Then for every $T>0$,
\[
P(N,T)\le8\|(1-\Pi)\delta_0\|^2+T^2\|\un_{B_N}\Pi\|^2.
\]
\end{lemma}

\begin{proof}
Let $\chi=\un_{[-N,N]}$. The commutator $[H,\chi]=[S+S^{-1},\chi]$ is supported on the two bonds $\{N,N+1\}$ and $\{-N-1,-N\}$, has norm one, and satisfies $[H,\chi]=[H,\chi]\un_{B_N}$. Since $\Pi$ commutes with $H$, Duhamel's formula gives
\[
(1-\chi)e^{-itH}\Pi\delta_0=e^{-itH}(1-\chi)\Pi\delta_0-i\int_0^te^{-i(t-s)H}[H,\chi]\,\Pi\,e^{-isH}\delta_0\,ds.
\]
As $(1-\chi)\delta_0=0$, we have $\|(1-\chi)\Pi\delta_0\|\le\|(1-\Pi)\delta_0\|$. Hence
\[
\|(1-\chi)\psi_t\|\le2\|(1-\Pi)\delta_0\|+|t|\,\|\un_{B_N}\Pi\|.
\]
Square, use $(x+y)^2\le2x^2+2y^2$, and average with $(2/T)e^{-2t/T}$, for which the mean of $t^2$ is $T^2/2$.
\end{proof}

The cutoff estimate contains a factor $T$ before squaring because Duhamel's formula integrates the boundary commutator over time. It controls the accumulated crossing error; by itself it does not give a time-uniform confinement bound.

\subsection{The local barrier at a scale}

We now apply these lemmas to $Q=P_q(H)$ from Proposition~\ref{prop:poly}, divided by $J_n$.

\begin{proposition}\label{prop:local}
Fix $c>0$ and $0<\eps<\min(L,c)/4$. There are $C$ and $n_0$ such that the following holds for every $n\ge n_0$ with $b_n\ge c$ and \emph{every} phase $\theta$. Write $s=|\sin(2\pi q_n\theta)|$, $\mu_n=\min(L,b_n)$ and $\kappa_n=-q_n^{-1}\log s\in[0,\infty]$. There is an integer $N_n$ with
\begin{equation}\label{eq:Nn}
e^{\eps q_n/2}\le N_n\le 3e^{\eps q_n}\min\Bigl\{\frac{\sigma_n}{|\delta_n|s},\ \frac{\sigma_n^{1/2}}{|\delta_n|}\Bigr\}\le 3\exp\Bigl(\bigl(b_n-\mu_n+\min\{\kappa_n,\mu_n/2\}+2\eps\bigr)q_n\Bigr)
\end{equation}
(the first entry of the minimum is $+\infty$ if $s=0$) and
\[
P(N_n,T)\le C(1+T^2)\exp\Bigl(-\frac{\eps N_n}{900}\Bigr)\quad\text{for all }T>0.
\]
In particular, if $s\ge e^{-\eps q_n}$, then $N_n\le3e^{((b_n-L)_++3\eps)q_n}$.
\end{proposition}

\begin{proof}
Write $q=q_n$, $\delta=\delta_n$, $b=b_n$, $A=\lambda^q$, $\mu=\mu_n\ge\mu_0=\min(L,c)$, $x=2\pi q\theta$ and $y_j=2\pi j\delta$, so that $V(j)=2A\cos(x+y_j)$. Take $\xi=\eps/8$ in Lemma~\ref{lem:sigma}, so that $|\delta|\le\sigma\le e^{-(\mu-\eps/4)q}\le1$, where $\sigma=J_n/A$. Put $s_*=e^{\eps q/2}\sigma^{1/2}$; since $\eps/2<(\mu-\eps/4)/2$, $s_*\to0$ uniformly in $b\ge c$. We construct $N$, an interval $I_0=[-r,r]$ with $r=\lfloor N/500\rfloor$, intervals $I_+=[N/2,N]$, $I_-=[-N,-N/2]$ and a number $V_0>0$ such that
\begin{equation}\label{eq:barriergeom}
\begin{gathered}
V(j)-V(0)\ \text{has a fixed sign on each of }I_+,\,I_-,\qquad |V(j)-V(0)|\ge V_0\ \text{on }I_+\cup I_-,\\
|V(j)-V(0)|\le V_0/32\ \text{on }I_0,\qquad V_0':=V_0/J_n\ge e^{\eps q/2},\qquad N\ge e^{\eps q/2}.
\end{gathered}
\end{equation}
The signs on $I_+$ and $I_-$ need not be opposite.

\emph{Case (i): $s\ge s_*$ (linear barrier).} Put $N=\lceil e^{\eps q/2}\sigma/(|\delta|s)\rceil$. Then $N\ge e^{\eps q/2}$ because $\sigma\ge|\delta|$ and $s\le1$, and
\[
2\pi|\delta|N\le2\pi\bigl(e^{\eps q/2}\sigma/s+|\delta|\bigr)\le2\pi(\sigma^{1/2}+\sigma)\le4\pi\sigma^{1/2}\le s_*/50\le s/50
\]
for large $q$. For $|j|\le N$ the mean value theorem gives $V(j)-V(0)=-2Ay_j\sin(x+\vartheta_jy_j)$ with $\vartheta_j\in(0,1)$, and $|\sin(x+\vartheta_jy_j)-\sin x|\le|y_j|\le s/50$. So $\sin(x+\vartheta_jy_j)$ has the sign of $\sin x$ and modulus in $[0.98s,1.02s]$. On $I_\pm$ we get a fixed sign and $|V(j)-V(0)|\ge2A\cdot\pi|\delta|N\cdot0.98s\ge A|\delta|Ns=:V_0$; on $I_0$, $|V(j)-V(0)|\le4\pi\cdot1.02\,A|\delta|rs\le V_0/32$. Finally $V_0'=A|\delta|Ns/J_n\ge e^{\eps q/2}A\sigma/J_n=e^{\eps q/2}$. Also $N\le2e^{\eps q/2}\sigma/(|\delta|s)$, and since $s\ge s_*$ this is at most $2\sigma^{1/2}/|\delta|$.

\emph{Case (ii): $s<s_*$ (quadratic barrier at a turning point).} Put $N=\lceil2e^{\eps q/2}\sigma^{1/2}/|\delta|\rceil$. Then $N\ge2e^{\eps q/2}$, since $\sigma^{1/2}\ge\sigma\ge|\delta|$. For $|j|\le N$,
\[
|y_j|\le2\pi|\delta|N\le2\pi(2s_*+|\delta|)\le0.01
\]
for large $q$, and $s<s_*\le0.01$, so $|\cos x|\ge0.99$. Taylor's formula gives
\[
V(j)-V(0)=-2Ay_j\sin x-Ay_j^2\cos(x+\vartheta_jy_j),\qquad\vartheta_j\in(0,1),
\]
and $|\cos(x+\vartheta_jy_j)-\cos x|\le|y_j|\le0.01$, so $\cos(x+\vartheta_jy_j)$ has the sign of $\cos x$ and modulus at least $0.98$. On $I_\pm$ we have $|y_j|\ge\pi|\delta|N\ge2\pi s_*>4s$, hence $2A|y_j|s\le\frac12Ay_j^2$, and therefore $V(j)-V(0)$ has the sign of $-\cos x$ on both $I_+$ and $I_-$, with
\[
|V(j)-V(0)|\ge0.48\,Ay_j^2\ge0.48\,A(\pi|\delta|N)^2=:V_0.
\]
On $I_0$, $|y_j|\le2\pi|\delta|N/500$, and $s<s_*\le|\delta|N/2$, so
\[
|V(j)-V(0)|\le2A|y_j|s+Ay_j^2\le A\pi^2\delta^2N^2\Bigl(\frac{2}{500\pi}+\frac{4}{500^2}\Bigr)\le0.0013\,A\pi^2\delta^2N^2\le V_0/32.
\]
Finally $V_0'=0.48\pi^2A\delta^2N^2/J_n\ge1.92\pi^2e^{\eps q}A\sigma/J_n\ge e^{\eps q/2}$. Also $N\le3e^{\eps q/2}\sigma^{1/2}/|\delta|$, and since $s<s_*$ we have $\sigma^{1/2}/|\delta|<e^{\eps q/2}\sigma/(|\delta|s)$, so $N\le3e^{\eps q}\sigma/(|\delta|s)$.

In both cases~\eqref{eq:barriergeom} holds and $N$ satisfies the first two inequalities in~\eqref{eq:Nn}. The third follows from $\sigma\le e^{-(\mu-\eps/4)q}$ and $|\delta|^{-1}=e^{bq}$: the two entries of the minimum are at most $e^{(b-\mu+\kappa_n+\eps/4)q}$ and $e^{(b-\mu/2+\eps/8)q}$.

\emph{The barrier estimates.} From here on the proof uses only~\eqref{eq:barriergeom}. Apply the lemmas to $Q'=Q/J_n=V'+K'$, $\|K'\|\le1$, with $d_0=V'(0)$ and $W'=V_0'/8$. Since $V_0'\ge e^{\eps q/2}$, $\log(7V_0'/8)\ge\eps q/3$ and $\log(3V_0'/32)\ge\eps q/3$ for large $q$, and $V_0'\ge W'+2$, $W'\ge4$. Let $R=\lfloor3N/4\rfloor$. Lemma~\ref{lem:outer}, applied separately on $I_+$ and on $I_-$ (on each of them $V'-d_0$ has a fixed sign and modulus at least $V_0'$), with $\ell=\lfloor(N/4-2)/q\rfloor-2\ge N/(5q)$, gives $\|\un_{B_R}\Pi\|\le4(8/(7V_0'))^{\ell}\le4e^{-\eps N/15}$. On $I_0$, $|V'-d_0|\le V_0'/32=W'/4$, and Lemma~\ref{lem:inner} with $\ell_0=\lfloor r/q\rfloor-2\ge N/(600q)$ gives $\|(1-\Pi)\delta_0\|\le3(32/(3V_0'))^{\ell_0}\le3e^{-\eps N/1800}$. Lemma~\ref{lem:cutoff} gives $P(R,T)\le72e^{-\eps N/900}+16T^2e^{-2\eps N/15}$, and $P(N,T)\le P(R,T)$. The last assertion of the proposition follows from~\eqref{eq:Nn} with $\kappa_n\le\eps$. All thresholds depend only on $\lambda,\alpha,c,\eps$, not on $\theta$. Take $N_n=N$.
\end{proof}

Near a turning point, case~(ii) replaces the linear barrier by a quadratic barrier that confines on both sides with the same sign. Its radius is of order $\sigma^{1/2}/|\delta|$, up to the stated exponential losses. For comparison, at a phase with sine $s_0>0$ fixed, the linear radius is of order $\sigma/(|\delta|s_0)$, and the ratio of these two scales is $s_0\sigma^{-1/2}$. Thus for $b_n>L$ the turning-point radius has exponential order $e^{(b_n-L/2)q_n}$, compared with $e^{(b_n-L)q_n}$ at such a non-turning phase. Remark~\ref{rem:tpnum} gives finite-volume numerical illustrations of these scales.

Two features matter later. The physical evolution is never rescaled; only $Q$ is divided by $J_n$, which changes neither its spectral projections nor its commutation with $H$. And the proposition needs only $b_n\ge c$: at scales with $b_n<L$ the available bound on $\|K\|$ may be large, but $\|K\|/\lambda^q$ is exponentially small.

\begin{remark}[Numerical illustration; not used in any proof]\label{rem:tpnum}
For $q=3$, $p=1$, $\alpha=(p+\delta)/q$ and $\lambda\in\{2,3,4\}$ we diagonalized $H_{\lambda,\alpha,\theta}$ on $[-M,M]$ with Dirichlet conditions and computed the infinite-time average $p_\infty(j)=\sum_E|\psi_E(0)|^2|\psi_E(j)|^2$ and its root mean square radius $r_{\mathrm{rms}}$; the mass within $20$ sites of the boundary was below $10^{-15}$ in all runs. At the turning points $\theta=0$ and $\theta=1/6$ (where $s=0$) we found $r_{\mathrm{rms}}\,|\delta|\lambda^{q/2}=0.136$--$0.143$ for all $\lambda\in\{2,3,4\}$ and $\delta\in\{10^{-3},2\cdot10^{-3},4\cdot10^{-3}\}$, that is, $r_{\mathrm{rms}}\approx0.14\,\lambda^{-q/2}/|\delta|$, the radius $\sigma^{1/2}/|\delta|$ of Proposition~\ref{prop:local}(ii). At $\theta=0.13$ ($s=0.637$) and $\lambda\in\{3,4\}$ we found $r_{\mathrm{rms}}\,|\delta|s\lambda^q=0.160$--$0.172$ (the largest value at $\lambda=4$, $\delta=4\cdot10^{-3}$, where the radius is about one site), the generic radius $\sigma/(|\delta|s)$; at $\lambda=4$, $\delta=10^{-3}$ the turning-point radius is $17.2$ against $3.9$ at $\theta=0.13$. The instantaneous radius at $\theta=0$ grows linearly, about $(1$--$2)\lambda^{-q}t$, until it saturates near $t\approx10^3$ for $\delta=10^{-3}$. In the regime $\frac12\log\lambda<b<\log\lambda$, where $K$ is not small and the generic exponential radius exponent $(b-L)_+$ is zero, and where Proposition~\ref{prop:local} at $s=0$ proves only the radius $e^{(b/2+O(\eps))q}$, which is larger than $\lambda^{-q/2}/|\delta|=e^{(b-L/2)q}$, we took $q=20$, $p=9$, $\lambda=2$ and $\delta\in\{6,3,1.5\}\cdot10^{-6}$, so $b=-\log|\delta|/q\in\{0.601,0.636,0.671\}<\log2=0.693$. At $\theta=0.013$ ($s=0.998$) we found $r_{\mathrm{rms}}=0.8$ in all three cases, while at $\theta=0$ we found $r_{\mathrm{rms}}=22.2,\,44.1,\,89.4$, against $0.14\,\lambda^{-q/2}/|\delta|=22.8,\,45.6,\,91.1$. These are single-scale finite-volume computations at rational-plus-perturbation frequencies; they say nothing directly about exponents.
\end{remark}
